\section{Integral products on compact threefolds}\label{sec:compact}

Let $\Gamma$ be an additive commutative monoid with an additive degree
$E:\Gamma\to\mathbb R_{\ge0}$, positive away from zero, with finite
sublevel sets. Multiplication in $H[[\Gamma]]$ is the convolution for
addition in $\Gamma$. An element has only finitely many decompositions
into nonzero summands: at its degree there are finitely many possible
sublabels and these have a positive minimum degree. For a series
$P=\sum P_\beta Q^\beta$ with $P_0=1$, compatibility with the
two geometric operations means
\begin{equation}\label{eq:twolaws}
 \Delta P_\beta=\sum_{\gamma+\delta=\beta}P_\gamma\otimes P_\delta,
 \qquad
 \rho_l P_\beta=\sum_{l\gamma=\beta}P_\gamma\quad(l>1).
\end{equation}
The second sum ranges over every root in $\Gamma$. Distinct integral
classes may have the same multiple.

\begin{proposition}\label{prop:criterion}
A compatible series has all its coefficients in $H_{\ZZ}$ if and only if it has
a product expansion in the factors $A_h(Q^\beta)$ with integral
exponents, finitely supported in $h$ at each $\beta\ne0$. In that case
the exponents are unique.
\end{proposition}
\begin{proof}
A single factor satisfies \eqref{eq:twolaws} by \eqref{eq:localops}.
Equivalently the second identity says that applying $\rho_l$
coefficientwise gives the substitution $Q^\gamma\mapsto Q^{l\gamma}$;
this identity and the coproduct identity persist under products and
inversion. All factors and their inverses have integral coefficients
by Theorem~\ref{thm:basis}.

Conversely divide out the factors already found at smaller degree.
At a first remaining degree the coefficient $R_\beta$ of the residual
integral series is primitive. Indeed in the first sum of
\eqref{eq:twolaws} every proper summand has smaller degree. For the
same reason every term on the right of the second identity vanishes.
Lemma~\ref{lem:detect} puts $R_\beta$ in $V$. By the integral
polynomial basis and $a_{h,1}=x_{h,1}$ we have
$V\cap H_{\ZZ}=\bigoplus_h\ZZ x_{h,1}$. The coefficients in that direct
sum are precisely the required exponents at $\beta$. Dividing out the
corresponding factors changes no other label of this degree.
Finite sublevel sets give the expansion and the same recursion proves
uniqueness.
\end{proof}

\begin{proof}[Proof of Theorem~\ref{thm:compact}]
On a compact Kähler threefold the mass of an effective cycle for the
Kähler metric is its area. Bounded-area cycles form a compact space
and their integral homology label is locally constant in the geometric
topology \cite[Theorem~2.11 and Remark~2.13]{DIW}. The image of a
compact space in this discrete group is finite, even when the group has
torsion. The cycles of fixed class give a compact open and closed
locus. A nonzero effective class has positive area. Nonnegativity of
$c_1(TX)$ implies that all components of a cycle in $\Gamma_0(X)$
have zero first Chern number. In particular the unit of the whole
locus is an integral semi-Fano insertion.

Addition of cycles pulls this unit back to the sum of the exterior
products of the units on the class loci for all splittings of $\beta$.
Multiplication by $l$ pulls it back to the units on all class loci
with $l\gamma=\beta$. These are precisely \eqref{eq:twolaws}, so
Proposition~\ref{prop:criterion} proves \eqref{eq:compact}. On a compact
manifold the areas of complex tangent lines under two Kähler metrics
are uniformly comparable. This identifies the two completions, and
uniqueness gives the same exponents. Under the projective hypotheses,
the geometric product of \cite[Theorem~1.1]{WueProduct} has exactly these
local factors and signed cluster multiplicities. Its rationalization
agrees with \eqref{eq:compact}; uniqueness identifies the integers.
\end{proof}
